Standard logit lenses decode intermediate language-model states into explicit tokenizer items, but recent evidence suggests that models also build word-like representations that need not correspond to a single token. We ask whether a lens can expose these implicit vocabulary items directly. We instantiate implicit tokens as 128 lowercased English word types from \wikitext that \qwen tokenizes into multiple pieces. For each layer, we build a layer-local implicit vocabulary table by averaging normalized hidden states from training occurrences, then decode held-out occurrences by cosine similarity to the word prototypes. On 512 held-out examples, the validation-selected \itl at layer 5 reaches 19.53\% top-1 accuracy, 25.59\% top-5 accuracy, and 0.240 MRR. This exceeds a final-subtoken majority baseline by 7.03 percentage points (bootstrap \ci [2.53, 11.72], paired randomization $p=0.0045$) and chance accuracy by 18.75 points. It also beats the strongest raw final-prototype implicit baseline by 5.47 points (bootstrap \ci [2.15, 8.79]). A Tuned-Lens-style ridge translator to final-layer prototypes performs poorly at this data scale, reaching only 3.71\% top-1. These results show that a supervised, prototype-based implicit-token lens can recover whole-word identity beyond subtoken leakage, while leaving open unsupervised discovery and causal validation of implicit vocabulary items.
